% ============ 第五章 §2 Abel L-函数, Чеботарёв 密度定理 (原书 pp.197-205) ============
\setcounter{example}{0}% 原书例编号分组

\section{\texorpdfstring{Abel $L$-函数, Чеботарёв 密度定理}{Abel L-函数, Chebotarev 密度定理}}\label{sec:ch5-s2}

设 $L/K$ 为数域的伽罗瓦扩张. 定理 \ref{thm:5.2} 表明, 在 $L$ 完全分裂的 $O_K$ 中素理想全体具有 $D$-密度 $1/[L:K]$. 现在我们要提出更一般的问题：对每对正整数 $f$ 和 $g$($fg=[L:K]$), 在 $L$ 上分解模式为 $(f,g)$ 的 $O_K$ 中素理想全体的 $D$-密度是多少？当 $f=1$, $g=[L:K]$ 时这就是完全分裂的情形. 

我们先对 Abel 扩张情形给出解答, 然后再考虑任意伽罗瓦扩张情形. 

\begin{theorem}\label{thm:5.3}
设 $L/K$ 是数域的 $n$ 次 Abel 扩张. 对 $G=\Gal(L/K)$ 中每个自同构 $\sigma$, 令
\[
  A(\sigma)=\left\{O_K\ \text{中素理想}\ \mathfrak{p}:\ \mathfrak{p}\ \text{在}\ L\ \text{上不分歧, 并且}\ \left(\frac{L/K}{\mathfrak{p}}\right)=\sigma\right\},
\]
则 $\delta(A(\sigma))=\frac{1}{n}$. 
\end{theorem}

\begin{proof}
映射
\[
  D_{\mathfrak{p}}\to\Gal(\overline{L}/\overline{K}),\qquad \sigma\mapsto\overline{\sigma}
\]
是群的满同态, 并且核为 $\mathfrak{p}$ 对 $L/K$ 的惯性群 $I_{\mathfrak{p}}$. 我们取 $\Gal(\overline{L}/\overline{K})$ 中自同构 $\tau\colon x\mapsto x^{|\overline{K}|}$(对 $x\in\overline{L}$)的一个原象 $\sigma_{\mathfrak{p}}\in D_{\mathfrak{p}}$, 则 $\tau$ 的全部原象为 $\sigma_{\mathfrak{p}}I_{\mathfrak{p}}$. 令 $\widehat{G}$ 为 Abel 群 $G=\Gal(L/K)$ 的特征群. 对每个特征 $\chi\in\widehat{G}$, 定义
\[
  \chi(\mathfrak{p})=\chi(\sigma_{\mathfrak{p}})\cdot\frac{1}{e(\mathfrak{P}/\mathfrak{p})}\sum_{\sigma\in I_{\mathfrak{p}}}\chi(\sigma)
  =\frac{1}{|I_{\mathfrak{p}}|}\sum_{\sigma\in I_{\mathfrak{p}}}\chi(\sigma_{\mathfrak{p}}\sigma).
\]
不难看出：

(I) $\chi(\mathfrak{p})$ 与 $\tau$ 的原象 $\sigma_{\mathfrak{p}}$ 之选取方式无关. 

(II) 若 $I_{\mathfrak{p}}$ 不包含在 $\operatorname{Ker}\chi$ 之中, 则 $\chi$ 在 $I_{\mathfrak{p}}$ 上不是平凡特征. 由特征的正交关系可知 $\chi(\mathfrak{p})=0$. 

(III) 若 $I_{\mathfrak{p}}\subseteq\operatorname{Ker}\chi$, 则 $\chi(\mathfrak{p})=\chi(\sigma_{\mathfrak{p}})$. 

(IV) 若 $\mathfrak{p}$ 在 $L$ 中不分歧, 则 $I_{\mathfrak{p}}=\{1\}$, 而 $\tau$ 只有一个原象 $\left(\frac{L/K}{\mathfrak{p}}\right)$, 从而 $\chi(\mathfrak{p})=\chi\left(\left(\frac{L/K}{\mathfrak{p}}\right)\right)$. 

现在将 $\chi$ 的定义完全积性地扩充到 $O_K$ 的全部非零整理想上, 即对于 $\mathfrak{a}=\mathfrak{p}_1\cdots\mathfrak{p}_r$, 定义 $\chi(\mathfrak{a})=\chi(\mathfrak{p}_1)\cdots\chi(\mathfrak{p}_r)$. 现在可定义一个 Dirichlet 级数
\begin{equation}\tag{1}
  L(s,\chi,L/K)=\sum_{\mathfrak{a}}\chi(\mathfrak{a})N(\mathfrak{a})^{-s},
\end{equation}
其中求和为 $\mathfrak{a}$ 过 $O_K$ 的全部非零整理想. 它叫作数域 Abel 扩张 $L/K$ 关于特征 $\chi\in\widehat{G}$ 的 \emph{Abel $L$-函数}. 这个函数有如下性质：

(a) 由于 $|\chi(\mathfrak{a})|=0$ 或 $1$, 可知级数 (1) 在 $\operatorname{Re}(s)>1$ 中定义出正则函数. 由于 $\chi$ 是完全积性的, 所以有 Euler 乘积公式
\[
  L(s,\chi,L/K)=\prod_{\mathfrak{p}}(1-\chi(\mathfrak{p})N(\mathfrak{p})^{-s})^{-1}\qquad(\operatorname{Re}(s)>1\ \text{时}).
\]

(b)
\[
  \prod_{\chi\in\widehat{G}}L(s,\chi,L/K)=\zeta_L(s)\qquad(\operatorname{Re}(s)>1),
\]
其中 $\zeta_L(s)$ 是 $L$ 的 Dedekind zeta 函数：
\[
  \zeta_L(s)=\prod_{\mathfrak{P}}(1-N(\mathfrak{P})^{-s})^{-1}\qquad(\operatorname{Re}(s)>1).
\]
由于 (b) 两边均有 Euler 展开, 所以只需对 $O_K$ 的每个素理想 $\mathfrak{p}$, 证明对应 $\mathfrak{p}$ 的因子相等：
\begin{equation}\tag{2}
  \prod_{\chi\in\widehat{G}}(1-\chi(\mathfrak{p})N(\mathfrak{p})^{-s})=\prod_{\mathfrak{B}\mid\mathfrak{p}}(1-N(\mathfrak{B})^{-s}).
\end{equation}
根据前述性质 (II) 和 (III), 可知 (2) 式左边为
\begin{equation}\tag{3}
  \prod_{\substack{\chi\in\widehat{G}\\ \operatorname{Ker}\chi\supseteq I_{\mathfrak{p}}}}(1-\chi(\mathfrak{p})N(\mathfrak{p})^{-s})
  =\prod_{\chi\in(G/I_{\mathfrak{p}})^{\wedge}}(1-\chi(\sigma_{\mathfrak{p}})N(\mathfrak{p})^{-s}).
\end{equation}
令 $\mathfrak{p}O_L=(\mathfrak{P}_1\cdots\mathfrak{P}_g)^e$, $efg=n$. 由于 $|G/I_{\mathfrak{p}}|=fg$, 而由 $\sigma_{\mathfrak{p}}$ 的取法知 $\sigma_{\mathfrak{p}}$ 在 $G/I_{\mathfrak{p}}$ 中阶为 $f$. 所以 (3) 式右边等于 $(1-N(\mathfrak{p})^{-fs})^g$. 另一方面, $\mathfrak{p}$ 的因子有 $g$ 个 $\mathfrak{P}_1,\cdots,\mathfrak{P}_g$, 而 $N(\mathfrak{P}_i)=\mathfrak{p}^f$. 于是 (2) 式右边也为 $(1-N(\mathfrak{p})^{-fs})^g$. 这就证明了 (b). 

(c) $L(s,\chi,L/K)$ 可以解析延拓到整个复平面上. 当 $\chi\neq\chi_0$ 时 $L(s,\chi,L/K)$ 在整个复平面上是正则的(没有奇点). 而 $L(s,\chi_0,L/K)=\zeta_K(s)$ 只在 $s=1$ 处有极点并且是单极点. 关于 (c) 的证明是不平凡的, 要用到类域论. 这里证明从略. 

有了以上准备, 现在可以很容易证明定理 \ref{thm:5.3}. 我们要证
\begin{equation}\tag{4}
  \lim_{s\to1^+}\sum_{\mathfrak{p}\in A(\sigma)}N(\mathfrak{p})^{-s}\big/-\log(s-1)=\frac{1}{n}.
\end{equation}
令
\[
  T(s)=\frac{1}{n}\sum_{\chi\in\widehat{G}}\chi(\sigma^{-1})\log L(s,\chi,L/K).
\]
由上面的 (b) 和 (c), 可知当 $\chi\neq\chi_0$ 时, $L(1,\chi,L/K)\neq 0,\infty$. 因此, 
\[
  \lim_{s\to1^+}T(s)/-\log(s-1)
  =\lim_{s\to1^+}\frac{1}{n}\cdot\frac{\log\zeta_K(s)}{-\log(s-1)}=\frac{1}{n}.
\]
但是
\begin{align*}
  T(s)&=\frac{1}{n}\sum_{\chi\in\widehat{G}}\chi(\sigma^{-1})\log\prod_{\mathfrak{p}}(1-\chi(\mathfrak{p})N(\mathfrak{p})^{-s})^{-1}\\
  &=\frac{1}{n}\sum_{\chi\in\widehat{G}}\chi(\sigma^{-1})\sum_{\mathfrak{p}}\log(1-\chi(\mathfrak{p})N(\mathfrak{p})^{-s})^{-1}\\
  &=\frac{1}{n}\sum_{\chi\in\widehat{G}}\chi(\sigma^{-1})\sum_{\mathfrak{p}}\sum_{m=1}^{\infty}\frac{\chi(\mathfrak{p})^mN(\mathfrak{p})^{-ms}}{m}\\
  &=\frac{1}{n}\sum_{\chi\in\widehat{G}}\chi(\sigma^{-1})\sum_{\mathfrak{p}}\chi(\mathfrak{p})N(\mathfrak{p})^{-s}+g(s)\\
  &=\frac{1}{n}\sum_{\mathfrak{p}}N(\mathfrak{p})^{-s}\sum_{\chi\in\widehat{G}}\chi(\sigma^{-1})\chi(\mathfrak{p})+g(s),
\end{align*}
其中 $g(s)$ 是对 $m\geqslant 2$ 求和部分, 如同以前一样可知 $g(s)$ 在 $s=1$ 附近是有界的. 而除了有限个 $\mathfrak{p}$ 之外, 均有 $\chi(\mathfrak{p})=\chi\left(\left(\frac{L/K}{\mathfrak{p}}\right)\right)$. 由特征正交关系知
\[
  \sum_{\chi\in\widehat{G}}\chi(\sigma^{-1})\chi\left(\left(\frac{L/K}{\mathfrak{p}}\right)\right)
  =\begin{cases}
    n, & \text{若}\ \sigma=\left(\frac{L/K}{\mathfrak{p}}\right),\\
    0, & \text{否则}.
  \end{cases}
\]
这就表明
\[
  T(s)=\sum_{\mathfrak{p}\in A(\sigma)}N(\mathfrak{p})^{-s}+g'(s),
\]
其中 $g'(s)$ 在 $s=1$ 附近有界. 综合上述可知
\[
  \frac{1}{n}=\lim_{s\to1^+}T(s)/-\log(s-1)
  =\lim_{s\to1^+}\sum_{\mathfrak{p}\in A(\sigma)}N(\mathfrak{p})^{-s}/-\log(s-1)
  =\delta(A(\sigma)).
\]
这就证明了定理 \ref{thm:5.3}. 
\end{proof}

\setcounter{corollary}{0}% 原书第五章系按节重新编号

\begin{corollary}[1]
设 $L/K$ 是 Abel 扩张, $fg=n=[L:K]$. 以 $n_f$ 表示 $G=\Gal(L/K)$ 中 $f$ 阶元素的个数, 而
\[
  A=\left\{O_K\ \text{中素理想}\ \mathfrak{p}:\ \mathfrak{p}\ \text{在}\ L\ \text{中素理想分解型式为}\ (f,g)\right\},
\]
则 $\delta(A)=\frac{n_f}{n}$. 
\end{corollary}

\begin{example}
设 $L/K$ 是数域的 Abel 扩张, 并且 $G=\Gal(L/K)$ 为两个 $2$ 阶循环群的直积, 则 $G$ 中有一个 $1$ 阶元素和 $3$ 个 $2$ 阶元素. 所以 $O_K$ 中满足 $\mathfrak{p}O_L=\mathfrak{P}_1\mathfrak{P}_2$ 的全部素理想 $\mathfrak{p}$ 其 $D$-密度为 $\frac{3}{4}$. 而在 $L$ 中完全分裂的全部 $\mathfrak{p}$ 具有 $D$-密度 $\frac{1}{4}$. 最后, 在 $L$ 中惯性的(即 $f=4$)全部 $\mathfrak{p}$ 具有 $D$-密度 $0$, 因为 $G$ 中没有 $4$ 阶元素. 注意：一个 $D$-密度为 $0$ 的集合不必为空集, 它可以是有限集合甚至是无限集合. 但事实上对于本例中的情形, 在 $L$ 中惯性的 $\mathfrak{p}$ 是不存在的. (为什么？)
\end{example}

\begin{remark}
(1) 由定理 \ref{thm:5.3} 的 (c) 可知
\[
  \zeta_L(s)/\zeta_K(s)=\prod_{\chi\neq\chi_0}L(s,\chi,L/K)
\]
是整个复平面上的正则函数. 换句话说, (利用类域论)可以证明 Artin 猜想对于 $L/K$ 为 Abel 扩张的情形是成立的. 

(2) 定理 \ref{thm:5.1} 是定理 \ref{thm:5.3} 的一个特例. 因为在定理 \ref{thm:5.3} 中取 $K=\Q$, $L=\Q(\zeta_k)$, 则 $[L:\Q]=\varphi(k)$, 而
\[
  G=\Gal(L/\Q)=\{\sigma_l:(l,k)=1,\ 1\leqslant l<k\},
\]
其中 $\sigma_l(\zeta_k)=\zeta_k^l$. 若 $p$ 在 $L$ 中不分歧(即 $p\nmid k$), 我们已经知道 $\left(\frac{L/K}{p}\right)=\sigma_p$, 所以 $\left(\frac{L/K}{p}\right)=\sigma_l\Leftrightarrow p\equiv l\ (\mathrm{mod}\ k)$. 所以素数集合 $A=\{p:p\equiv l\ (\mathrm{mod}\ k)\}$ 的 $D$-密度就是素数集合 $A'=\left\{p:\left(\frac{L/K}{p}\right)=\sigma_l\right\}$ 的 $D$-密度. 根据定理 \ref{thm:5.3}, $\delta(A')=\frac{1}{\varphi(k)}$, 从而 $\delta(A)=\frac{1}{\varphi(k)}$. 

(3) 我们略去了定理 \ref{thm:5.2} 的一部分证明, 因为证明 Abel $L$-函数解析延拓和奇点特性时要用到类域论. 但是当 $K=\Q$ 时, 我们可以不用类域论而只用 Kronecker-Weber 定理证明定理 \ref{thm:5.3}. 因为若 $L/\Q$ 是 Abel 扩张, 则 $L$ 是分圆域 $\Q(\zeta_m)$ 的子域, 其中 $m$ 为域 $L$ 的导子. 伽罗瓦群 $\Gal(\Q(\zeta_m)/\Q)$ 自然同构于 $(\Z/m\Z)^{\times}$. 设 $N$ 为 $L$ 的固定子群, $N$ 在 $(\Z/m\Z)^{\times}$ 中的象为 $\varphi(N)$, 则 $G=\Gal(L/\Q)$ 自然同构于 $(\Z/m\Z)^{\times}/\varphi(N)$(参见定理 \ref{thm:2.18} 的证明和后面的注记 1). 若素数 $p$ 在 $L$ 中不分歧, 我们不妨设 $(p,m)=1$(因为 $m$ 的有限个素因子是 $D$-密度为 $0$ 的集合, 不影响以后的讨论). $G$ 中每个元素 $\sigma$(即 $L$ 的自同构)自然看成是 $(\Z/m\Z)^{\times}$ 模 $\varphi(N)$ 的一个陪集, 所以相当于模 $m$ 的 $|\varphi(N)|$ 个同余类. 注意 $|\varphi(N)|=|N|=[\Q(\zeta_m):L]=\frac{\varphi(m)}{[L:\Q]}$. 如果 $\left(\frac{L/\Q}{p}\right)=\sigma_{\mathfrak{p}}$, 所以这相当于 $p$ 属于 $\sigma$ 所对应的那 $\frac{\varphi(m)}{[L:\Q]}$ 个模 $m$ 同余类. 根据例 1, 每个同余类中素数的 $D$-密度为 $\frac{1}{\varphi(m)}$. 所以满足 $\left(\frac{L/\Q}{p}\right)=\sigma$ 的所有素数组成的集合, 其 $D$-密度为 $\frac{\varphi(m)}{[L:\Q]}\cdot\varphi(m)^{-1}=[L:\Q]^{-1}$. 这就是定理 \ref{thm:5.3}(对于 $K=\Q$ 的情形). 
\end{remark}

现在讨论 $L/K$ 为伽罗瓦扩张的情形. 这时, 设 $O_K$ 中素理想 $\mathfrak{p}$ 在 $L$ 中不分歧, 则 $\mathfrak{p}O_L=\mathfrak{P}_1\cdots\mathfrak{P}_g$, 令 $\mathfrak{P}=\mathfrak{P}_1$, 则每个 $\mathfrak{P}_i$ 均可写成 $\sigma(\mathfrak{P})$(对某个 $\sigma\in G=\Gal(L/K)$). 但是
\[
  \left(\frac{L/K}{\sigma(\mathfrak{P})}\right)=\sigma\left(\frac{L/K}{\mathfrak{P}}\right)\sigma^{-1},
\]
于是 $\left\{\left(\frac{L/K}{\mathfrak{P}_i}\right)\mid 1\leqslant i\leqslant g\right\}=\left\{\sigma\left(\frac{L/K}{\mathfrak{P}}\right)\sigma^{-1}\mid\sigma\in G\right\}$, 这是有限群 $G$ 的一个共轭类, 我们将这个共轭类表示成 $\left(\frac{L/K}{\mathfrak{p}}\right)$(当 $L/K$ 为 Abel 扩张时, 这个共轭类只含有一个元素, 我们曾经把这个元素表示成 $\left(\frac{L/K}{\mathfrak{p}}\right)$). 

\begin{theorem}[\textbf{Чеботарёв}]\label{thm:5.4}
设 $L/K$ 是数域的伽罗瓦扩张, $G=\Gal(L/K)$. $C$ 为 $G$ 中一个共轭元素类, $|C|=c$. $A=\left\{O_K\ \text{中素理想}\ \mathfrak{p}\mid\mathfrak{p}\ \text{在}\ L\ \text{中不分歧并且}\ \left(\frac{L/K}{\mathfrak{p}}\right)=C\right\}$, 则 $\delta(A)=c/n$, $n=[L:K]$. 
\end{theorem}

\begin{proof}
俄国数学家 Чеботарёв 原来的证明很复杂. 这里采用 $1969$ 年 McCluer 给出的简化证明. 首先需要下列引理. 

\setcounter{lemma}{4}% 原书第五章引理编号跳号(无引理2-4)

\begin{lemma}\label{lem:ch5-5}
设 $L/K$ 是数域的伽罗瓦扩张, $\mathfrak{P}$ 和 $\mathfrak{p}$ 分别是 $O_L$ 和 $O_K$ 中的(固定)素理想, $\mathfrak{P}\mid\mathfrak{p}$, $\mathfrak{p}$ 在 $L$ 中不分歧. $M$ 是 $L/K$ 的中间域. 并且对于 $O_M$ 中每个素理想 $\mathfrak{q}\mid\mathfrak{p}$, $\mathfrak{q}$ 在 $L$ 中均是惯性的(从而 $\mathfrak{q}O_L$ 为 $O_L$ 中素理想). 令 $\left(\frac{L/K}{\mathfrak{P}}\right)=\sigma\in G=\Gal(L/K)$, 则 $\mathfrak{p}$ 在 $O_M$ 中共存在 $[C_G(\sigma):H]$ 个素理想因子 $\mathfrak{q}$, 使得 $\left(\frac{L/K}{\mathfrak{q}O_L}\right)=\sigma$. 其中 $C_G(\sigma)=\{\tau\in G\mid\tau\sigma=\sigma\tau\}$($\sigma$ 在 $G$ 中的中心化子), 而 $H=\langle\sigma\rangle$($\sigma$ 在 $G$ 中生成的循环子群). 
\end{lemma}

这是因为：从引理的假设可知, 我们只需证明 $O_L$ 中共有 $[C_G(\sigma):H]$ 个 $\mathfrak{P}'\mid\mathfrak{p}$ 使得 $\left(\frac{L/K}{\mathfrak{P}'}\right)=\sigma$. 每个这种 $\mathfrak{P}'$ 均可写成 $\tau(\mathfrak{P})$, $\tau\in G$. 而
\[
  \left(\frac{L/K}{\tau(\mathfrak{P})}\right)=\sigma
  \Leftrightarrow\tau\sigma\tau^{-1}=\sigma
  \Leftrightarrow\tau\in C_G(\sigma);
\]
并且
\[
  \tau_1\mathfrak{P}=\tau_2\mathfrak{P}
  \Leftrightarrow\tau_2^{-1}\tau_1\mathfrak{P}=\mathfrak{P}
  \Leftrightarrow\tau_2^{-1}\tau_1\in D_{\mathfrak{P}}=\langle\sigma\rangle=H.
\]
由此即得引理. 

现在证明定理 \ref{thm:5.4}：这只需证明
\[
  \sum_{\left(\frac{L/K}{\mathfrak{p}}\right)=C}N(\mathfrak{p})^{-s}
  =-\frac{c}{n}\log(s-1)+o(\log(s-1))\qquad(\text{当 } s\to1^+\ \text{时}).
\]
取 $\sigma\in C$, $H=\langle\sigma\rangle\subseteq G=\Gal(L/K)$. $M$ 为 $H$ 的固定子域, 则 $\Gal(L/M)\cong H$, 从而 $L/M$ 为循环扩张. 由定理 \ref{thm:5.3} 知
\[
  \sum_{\left(\frac{L/M}{\mathfrak{q}}\right)=\sigma}N(\mathfrak{q})^{-s}
  =-\frac{1}{|H|}\log(s-1)+o(\log(s-1)).
\]
从而
\[
  \sum_{\substack{\mathfrak{q}\mid\mathfrak{p}\\ \left(\frac{L/M}{\mathfrak{q}}\right)=\sigma\\ f(\mathfrak{q}\mid\mathfrak{p})=1}}N(\mathfrak{p})^{-s}
  =-\frac{1}{|H|}\log(s-1)+o(\log(s-1)).
\]
对于使 $\left(\frac{L/M}{\mathfrak{q}}\right)=\sigma$ 和 $\mathfrak{q}\mid\mathfrak{p}$ 成立的每个 $\mathfrak{p}$, $M$ 满足上面引理条件. 因为 $\left(\frac{L/M}{\mathfrak{q}}\right)=\sigma$, 则分解群 $D_{\mathfrak{p}}=\langle\sigma\rangle\cong\Gal(L/M)$, 从而 $\mathfrak{q}$ 在 $L$ 中惯性. 若 $\mathfrak{q}'$ 为 $M$ 中素理想并且 $\mathfrak{q}'\mid\mathfrak{p}$, 则 $D_{\mathfrak{q}'}$ 与 $D_{\mathfrak{q}}$ 共轭, 从而仍有 $|D_{\mathfrak{q}'}|=|\Gal(L/M)|$, 则 $\mathfrak{q}'$ 在 $L$ 中惯性. 于是由引理得出(注意 $\left(\frac{L/M}{\mathfrak{q}}\right)=\sigma\Leftrightarrow\left(\frac{L/K}{\mathfrak{p}}\right)=C$)
\[
  [C_G(\sigma):H]\sum_{\left(\frac{L/K}{\mathfrak{p}}\right)=C}N(\mathfrak{p})^{-s}
  =-\frac{1}{|H|}\log(s-1)+o(\log(s-1)).
\]
于是
\[
  \sum_{\left(\frac{L/K}{\mathfrak{p}}\right)=C}N(\mathfrak{p})^{-s}
  =-\frac{1}{[C_G(\sigma):H]\cdot|H|}\log(s-1)+o(\log(s-1)).
\]
但是 $[C_G(\sigma):H]\cdot|H|=|C_G(\sigma)|=\frac{n}{c}$, 这就表明 $\delta(A)=\frac{c}{n}$. 
\end{proof}

我们可以把定理 \ref{thm:5.3} 的系 1 推广到任意伽罗瓦扩张的情形. 

\begin{corollarynonum}
设 $L/K$ 是数域的伽罗瓦扩张, $fg=n=[L:K]$. 以 $n_f$ 表示 $G=\Gal(L/K)$ 中 $f$ 阶元素的个数, 令
\[
  A=\{O_K\ \text{的素理想}\ \mathfrak{p}:\mathfrak{p}\ \text{在}\ L\ \text{中分解型式为}\ (f,g)\},
\]
则 $\delta(A)=\frac{n_f}{n}$. 
\end{corollarynonum}

\begin{proof}
设 $\mathfrak{p}$ 在 $L$ 中不分歧, 则 $\mathfrak{p}$ 在 $L$ 中分解型式为 $(f,g)$ 当且仅当 $G$ 的共轭元素类 $\left(\frac{L/K}{\mathfrak{p}}\right)$ 中所有元素的阶均为 $f$. 再由定理 \ref{thm:5.4} 即得此系. 
\end{proof}

% ============ 习题 ============
\begin{Exercises}
  \item 设 $f(x)$ 和 $g(x)$ 为 $\Z[x]$ 中不可约多项式, $\alpha$ 和 $\beta$ 分别是 $f(x)$ 和 $g(x)$ 在复数域中的一个根. 如果 $\Q(\alpha)\subseteq\Q(\beta)$, 求证 $S(f)$ 几乎包含 $S(g)$, 即素数集合 $\{p:p\in S(g),\ p\notin S(f)\}$ 的 $D$-密度为 $0$. 
  \item 设 $f_1(x),\cdots,f_r(x)$ 是 $\Z[x]$ 中有限个多项式, 并且次数均 $\geqslant 1$. 求证存在无限多个素数 $p$, 使得 $f_1(x),\cdots,f_r(x)$ 模 $p$ 均完全分裂. 
  \item 设 $L/K$ 是数域的伽罗瓦扩张. 求证：$O_K$ 中有无限多个素理想 $\mathfrak{p}$ 在 $L$ 中惯性的充分必要条件为 $L/K$ 是循环扩张(即伽罗瓦群 $\Gal(L/K)$ 为循环群). 
  \item 计算素数集合 $A=\{p:2$ 为模 $p$ 的三次剩余$\}$ 和 $A'=\{p:2$ 为模 $p$ 的四次剩余$\}$ 的 $D$-密度. 
  \item 设 $f(x)$ 是 $\Z[x]$ 中次数 $\geqslant 2$ 的首 $1$ 不可约多项式, 求证存在无限多素数 $p$, 使得 $f(x)$ 在有限域 $\F_p$ 中无根. 
\end{Exercises}
